\section{第\ref{sec:eetypes}章　素子としてのニューロン}

個々の細胞の動作モードは、電極からの電流注入と電圧記録で直接測定されている。積分器とクラス分けの数学は古典的な理論である。脳がモードの切り替えを計算に使っているというのは、依然として仮説である。

{\footnotesize
\setlength{\tabcolsep}{3pt}\renewcommand{\arraystretch}{1.15}
\begin{longtable}{>{\raggedright}p{0.5cm}>{\raggedright}p{4.0cm}>{\raggedright}p{5.6cm}>{\raggedright}p{2.3cm}>{\raggedright\arraybackslash}p{1.7cm}}
\toprule
№ & 主張 & 実験と測定内容 & 出典 & 状態 \\
\midrule\endfirsthead
\toprule
№ & 主張 & 実験と測定内容 & 出典 & 状態 \\
\midrule\endhead
1 & 共振器：電圧の変化に逆らう遅い電流が、ニューロンを数Hzから数十Hzに最大をもつバンドパスフィルタにする（第\ref{sec:resonator}節） & スライス中の\nt{GLUT}ニューロンに周波数が滑らかに上がる正弦波電流を注入し、インピーダンス$|Z(f)|$を測定：最大は$33^\circ$Cで$2$--$5$~Hz（$38^\circ$Cで約$7$~Hz）。共振は遅いカリウム電流と陽イオン電流が生み、持続性ナトリウム電流が強める。総説は多くの細胞クラスで同じ手法を示す & \cite{HuVervaekeStorm2002,HutcheonYarom2000} & 測定済み \\
2 & 遅い電流はインダクタンスのように振る舞い、膜容量と合わせて共振回路を作る & 電流に対する軸索の閾値下応答を測定し、線形化したコンダクタンス方程式で記述。等価回路は電圧に依存する「現象論的」インダクタンスを含む & \cite{Mauro1970} & 理論 \\
3 & フィードバックのある群の時定数は$\tau_{\text{eff}}=\tau/(1-w)$~\eqref{eq:perfectint}。$\tau=0.1$~s、$w=0.995$で$20$~s & 群の線形方程式から本章で導出。眼の位置を保持する機構として理論研究で提案された & 本章で導出；\cite{Seung1996} & 理論 \\
4 & 眼の位置の積分器は、個々の細胞の性質ではなくネットワークで入力を保持する & 眼の位置を保持する魚の神経核での細胞内記録：急速な回転のあと、ニューロンの発火率は階段状に変わって保たれる。1個の細胞への短い電流は一時的なずれしか起こさず、電圧の階段は閾値下でも保たれる。それを維持するのはシナプス入力である & \cite{Aksay2001} & 測定済み \\
5 & 漏れのない積分器には学習による絶え間ない調整が必要。調整が狂うと忘れるか飽和する & 眼の位置の$\pm$に比例する視覚フィードバックで魚を訓練：保持は不安定または漏れのあるものになり、ニューロンの時定数は1桁以上下がって$<1$~sになった。通常の視覚フィードバックで安定性は徐々に戻った & \cite{Major2004} & 測定済み \\
6 & 双安定ニューロン：短い入力が長いプラトーをオンにし、抑制がオフにする。この性質はセロトニンでオンになる（第\ref{sec:bistable}節） & ネコの筋肉を制御する\nt{ACET}ニューロンの細胞内記録：短いシナプス興奮でニューロンは長く続く発火に移り、短い抑制でリセットされる。脊髄ネコではセロトニン前駆体の投与後にこの状態が現れる & \cite{Hounsgaard1988} & 測定済み \\
7 & 二つのクラス：積分器型（発火率がゼロから上がる）と共振器型（閾値ですぐ有限の発火率） & クラスは分岐理論から導かれる。実験：皮質スライスで\nt{GLUT}ニューロンはいくらでも低い発火率で発火し始め（タイプ1）、速い\nt{GABA}ニューロンは有限の発火率から跳躍的に始まる（タイプ2） & \cite{Izhikevich2007,Tateno2004} & 測定済み \\
8 & 1個のニューロンが積分器にも同時検出器にもなる：抑制が加算の窓を短くする & 海馬スライスで、興奮性入力が閾値まで加算される窓は$2$~ms未満。興奮の直後に届くフィードフォワード抑制がこれを短くする。抑制の弱い樹状突起では窓が広い & \cite{Pouille2001} & 測定済み \\
9 & 軸索による遅延線（表~\ref{tab:eetypes}） & メンフクロウで、同時検出器へ向かう軸索の遅延を測定：軸索の長さの違いが異なる遅延を生み、検出器は音の到達時間差に同調している & \cite{Carr1990} & 測定済み \\
10 & 覚醒中はペースメーカー、睡眠中は黙った細胞 & \nt{SERO}ニューロンは日中ほぼ規則的に発火し、徐波睡眠で遅くなり、レム睡眠ではほぼ沈黙する（詳しくは第\ref{sec:sleep}章） & \cite{JacobsAzmitia1992,McGintyHarper1976} & 測定済み \\
11 & 素子を決めるのはイオンチャネルと、それを調整するブロードキャストチャネル（命題~\ref{prop:eetypes}） & 小さなネットワークで示された：調節物質がニューロン固有の性質とシナプスの強さを変え、同じ回路が異なるリズムを出す & \cite{Marder2012} & 測定済み \\
12 & 脳はモードの切り替えを計算に使っている（命題~\ref{prop:eetypes}） & 細胞のモードの切り替えが課題の解決に必要であることを示す直接の実験はない。調節物質が回路の出力を変える実験があるだけである & \cite{Marder2012} & 仮説 \\
\bottomrule
\end{longtable}}
